\documentclass[../../../include/open-logic-section]{subfiles}

\begin{document}

\olfileid[hi]{sth}{z}{pairs}
\olsection{युग्म}

विचार करने योग्य अगला अभिगृहीत यह है:

\begin{axiom}[युग्म]
किन्हीं समुच्चयों $a, b$ के लिए समुच्चय $\{a, b\}$ का अस्तित्व है।
\[
	\forall a \forall b \exists P \forall x (x \in P \liff (x = a \lor x = b))
\]
\end{axiom}

पुनरावर्ती अवधारणा से इस अभिगृहीत का औचित्य इस प्रकार दिया जा सकता है।
मान लें कि $a$ चरण $S$ पर और $b$ चरण $T$ पर उपलब्ध है। $M$ को $S$ और
$T$ में से जो बाद में आता है वह चरण मानें। चूँकि $a$ और $b$ दोनों चरण $M$
पर उपलब्ध हैं, समुच्चय $\{a,b\}$ एक संभव संग्रह है जो $M$ के बाद किसी भी
चरण पर उपलब्ध है (दोनों में जो बड़ा हो)।

पर रुकिए!{} यह क्यों मानें कि $M$ के बाद कोई चरण \emph{है}? यदि कोई नहीं,
तो हमारा औचित्य विफल होगा। अतः युग्म का औचित्य देने के लिए
\olref[sth][z][story]{sec} की कथा में एक और सिद्धांत जोड़ना होगा:
\begin{enumerate}
	\item[] \stagessucc. कोई अंतिम चरण नहीं है।
\end{enumerate}
क्या यह सिद्धांत उचित है? शोएनफ़ील्ड की कथा में कहीं भी \emph{स्पष्टतः}
नहीं कहा गया कि कोई अंतिम चरण नहीं है। फिर भी यदि वह (ठीक-ठीक कहें तो)
हमारी कथा में एक अतिरिक्त जोड़ है, तो भी वह इस मूल विचार के अनुकूल है कि
समुच्चय चरणों में बनते हैं। आगे हम इसे बस स्वीकार करेंगे। और इसलिए युग्म का
अभिगृहीत भी स्वीकार करेंगे।

इस नए अभिगृहीत से हम और बहुत-से समुच्चयों का अस्तित्व सिद्ध कर सकते हैं।
उदाहरणतः:

\begin{prop}\ollabel{prop:pairsconsequences}
किन्हीं समुच्चयों $a$ और $b$ के लिए निम्न समुच्चयों का अस्तित्व है:
	\begin{enumerate}
		\item\ollabel{singleton} $\{a\}$
		\item\ollabel{binunion} $a \cup b$
		\item\ollabel{tuples} $\tuple{a, b}$
	\end{enumerate}
\end{prop}

\begin{proof}
\olref{singleton}. युग्म से $\{a, a\}$ का अस्तित्व है, जो विस्तारिता से
$\{a\}$ है।

\olref{binunion}. युग्म से $\{a, b\}$ का अस्तित्व है। अब सम्मिलन से
$a \cup b = \bigcup \{a, b\}$ का अस्तित्व है।

\olref{tuples}. \olref{singleton} से $\{a\}$ का अस्तित्व है। युग्म से
$\{a, b\}$ का अस्तित्व है। अब फिर युग्म से
$\{\{a\}, \{a, b\}\} = \tuple{a, b}$ का अस्तित्व है।
\end{proof}

\begin{prob}
दिखाइए कि किन्हीं समुच्चयों $a, b, c$ के लिए समुच्चय $\{a, b, c\}$ का
अस्तित्व है।
\end{prob}

\begin{prob}
दिखाइए कि किन्हीं समुच्चयों $a_1, \ldots, a_n$ के लिए समुच्चय
$\{a_1, \ldots, a_n\}$ का अस्तित्व है।
\end{prob}

%\begin{proof} By two applications of Pairs, $\{\{a_1, a_2\}, \{a_1,
%   a_3\}\}$ exists. By Union and Extensionality, $\bigcup \{\{a_1,
%   a_2\}, \{a_1, a_3\}\} = \{a_1, a_2, a_3\}$ exists. Repeat this
%   trick as often as necessary. \end{proof}

\end{document}
